\section*{第 \ref{ch:b2:liquid-vapour-diagrams} 章——二元液–蒸气相图与蒸馏}

\begin{solution}{exo:b2:liquid-vapour-diagrams:1}
苯 \qty{80.1}{\degreeCelsius}，甲苯 \qty{110.6}{\degreeCelsius}。
组成为 0.2 的液体约在 \qty{102}{\degreeCelsius} 开始沸腾，其
首批蒸气的组成为 0.38。
\end{solution}

\begin{solution}{exo:b2:liquid-vapour-diagrams:2}
$p = 0.5 \times 1.010 + 0.5 \times 0.388 = \qty{0.699}{bar}$;
$y = 0.505/0.699 = 0.722$.
\end{solution}

\begin{solution}{exo:b2:liquid-vapour-diagrams:3}
$n_V/n = (0.30 - 0.20)/(0.45 - 0.20) = 0.40$：\qty{4.0}{mol} 蒸气和
\qty{6.0}{mol} 液体。核验：$6.0 \times 0.20 + 4.0 \times 0.45 = 3.0 = 10
\times 0.30$.
\end{solution}

\begin{solution}{exo:b2:liquid-vapour-diagrams:4}
(a) $v = 2 + 2 - 1 = 3$：$T$、$p$ 和 $x$ 可独立变化。(b) $v = 2 + 2 - 2 = 2$；在
固定压强时为 1。(c) \omterm{def:b2:liquid-vapour-diagrams:azeotrope}{共沸物}增加了关系 $x = y$：$v = 1$；在固定
压强下为 0——\omterm{def:b2:liquid-vapour-diagrams:azeotrope}{共沸物}与纯液体一样在固定温度下沸腾，
但其组成随压强而变。
\end{solution}

\begin{solution}{exo:b2:liquid-vapour-diagrams:5}
在\omterm{def:b2:liquid-vapour-diagrams:bubble-dew}{露点}处 $yp/p_1^* + (1 - y)p/p_2^* = 1$。在 \qty{100}{\degreeCelsius} 时：
$0.3 \times 1.013/1.80 + 0.7 \times 1.013/0.742 = 1.124 > 1$（温度过低）；在
\qty{105}{\degreeCelsius} 时：$0.148 + 0.824 = 0.971 < 1$。内插得 $T \approx
100 + 5 \times 0.124/0.153 = \qty{104}{\degreeCelsius}$。首滴液体的 $x =
yp/p_1^* \approx 0.3 \times 1.013/2.0 = 0.15$.
\end{solution}

\begin{solution}{exo:b2:liquid-vapour-diagrams:6}
首批液滴的组成与沸腾液体上方蒸气的组成相同，
即 $y = 0.71$，温度为 \qty{92.1}{\degreeCelsius}。由于苯优先逸出，
液体中苯的含量降低，沸点升高，馏出物
中苯的含量同样降低；最后的液滴几乎是纯甲苯。简单蒸馏相当于一次
平衡级，只是在不断变化的液体上重复进行：每一
馏分都是混合物。
\end{solution}

\begin{solution}{exo:b2:liquid-vapour-diagrams:7}
进料组成 0.05，位于\omterm{def:b2:liquid-vapour-diagrams:azeotrope}{共沸物}的水一侧：塔顶得到\omterm{def:b2:liquid-vapour-diagrams:azeotrope}{共沸物}
（按摩尔计为 0.88，按质量计为 95~\%），塔釜保留水。进料组成为 0.95，位于
乙醇一侧：塔顶仍得到\omterm{def:b2:liquid-vapour-diagrams:azeotrope}{共沸物}，其沸点最低，
塔釜保留纯乙醇。
\end{solution}

\begin{solution}{exo:b2:liquid-vapour-diagrams:8}
混合物在 $0.10 + p^*_{\text{水}} = 1.013$ 时沸腾，即
$p^*_{\text{水}} = \qty{0.91}{bar}$，接近 \qty{97}{\degreeCelsius}。质量比
$0.10 \times 136/(0.91 \times 18.02) = 0.83$：每克
水对应 \qty{0.83}{g} 精油。
\end{solution}

\begin{solution}{exo:b2:liquid-vapour-diagrams:9}
室温下存在两层液体 L$_2$ 和 L$_1$（组成 0.4 位于
\omterm{def:b2:liquid-vapour-diagrams:miscibility-gap}{互溶间隙}内）。加热时，两层液体在\omterm{def:b2:liquid-vapour-diagrams:miscibility-gap}{非均相共沸物}的
温度下一起沸腾，且该温度保持不变；首批蒸气的组成为非均相共沸组成
（图中的点）。由于 0.4 位于该点的 L$_2$ 一侧，液层
L$_1$ 先被耗尽；随后温度升高，剩余的 L$_2$
沿其\omterm{def:b2:liquid-vapour-diagrams:bubble-dew}{泡点线}组分 1 含量下降，蒸气沿\omterm{def:b2:liquid-vapour-diagrams:bubble-dew}{露点线}变化，直至
该点到达\omterm{def:b2:liquid-vapour-diagrams:bubble-dew}{露点线}时，最后一滴液体蒸发。
\end{solution}

\begin{solution}{exo:b2:liquid-vapour-diagrams:10}
$N_{\min} = \ln(99 \times 99)/\ln 2.47 = 9.19/0.904 = 10.2$，对比 6.5：
两端纯度由 95~\% 提高到 99~\% 时，所需塔板数
再增加超过一半。$x/(1 - x)$ 上每增加一个因子，所需塔板数相同，而
纯度以残留杂质衡量，它只按几何级数下降。
\end{solution}

\begin{solution}{exo:b2:liquid-vapour-diagrams:11}
仅由\omterm{def:b2:liquid-vapour-diagrams:binary-diagram}{等压相图}无法得出，需要 $p_1^*$ 与 $p_2^*$ 在
\qty{80}{\degreeCelsius} 时的值，由安托万方程求得（1.010 与
\qty{0.388}{bar}）。于是 $p = p_2^* + x(p_1^* - p_2^*)$ 和 $p = 1/(y/p_1^* + (1 -
y)/p_2^*)$。液体在高压下稳定（压缩蒸气可使其
液化），在\omterm{def:b2:liquid-vapour-diagrams:binary-diagram}{等压相图}上则位于低温一侧：两张图
彼此上下颠倒。
\end{solution}

\begin{solution}{exo:b2:liquid-vapour-diagrams:12}
由于 $y$ 随 $x$ 增大且 $y(1 - y) > 0$，$dp/dx$ 的符号与 $y - x$ 相同：
在蒸气中组分 1 较富集处，压强随 $x$ 增大。对\omterm{def:b2:chemical-potential:ideal-mixture}{理想混合物}，
$y > x$ 对每个 $0 < x < 1$ 成立（见\cref{prop:b2:liquid-vapour-diagrams:richer}）：
$p$ 严格单调，无极值，也不存在\omterm{def:b2:liquid-vapour-diagrams:azeotrope}{共沸物}。
\end{solution}

\begin{solution}{pb:b2:liquid-vapour-diagrams:1}
\textbf{1.} $T = B/(A - \log_{10}1.013) - C$：苯 $1203.835/4.01253 + 53.226 =
\qty{353.2}{K}$，\qty{80.1}{\degreeCelsius}；甲苯 $1343.943/4.07266 + 53.773 =
\qty{383.8}{K}$，\qty{110.6}{\degreeCelsius}。
\textbf{2.} 在 \qty{363.15}{K} 时：$p_1^* = \qty{1.361}{bar}$，$p_2^* = \qty{0.542}{bar}$。
\textbf{3.} $x = (1.013 - 0.542)/(1.361 - 0.542) = 0.575$；$y = 0.575 \times
1.361/1.013 = 0.773$.
\textbf{4.} $x = (1.013 - 0.742)/(1.800 - 0.742) = 0.256$；$y = 0.256 \times
1.800/1.013 = 0.455$.
\textbf{5.} \omterm{def:b2:liquid-vapour-diagrams:bubble-dew}{泡点}：(0, 110.6)，(0.256, 100)，(0.575, 90)，(1, 80.1)；\omterm{def:b2:liquid-vapour-diagrams:bubble-dew}{露
点}：(0, 110.6)，(0.455, 100)，(0.773, 90)，(1, 80.1)——图中的
透镜形区域。
\textbf{6.} $\frac12(p_1^* + p_2^*) = 1.013$：在 \qty{90}{\degreeCelsius} 时为 0.952，在
\qty{95}{\degreeCelsius} 时为 1.102；内插得 $90 + 5 \times 0.061/0.150 \approx
\qty{92}{\degreeCelsius}$（恰为 \qty{92.1}{\degreeCelsius}）。
\textbf{7.} $v = 2$，固定压强时为 1：每个温度都确定 $x$ 和 $y$，
这正是该图由两条曲线构成的原因。
\textbf{8.} $x = (1.013 - 0.636)/(1.569 - 0.636) = 0.404$；$y = 0.404 \times
1.569/1.013 = 0.626$.
\textbf{9.} $n_V = 100 \times (0.500 - 0.404)/(0.626 - 0.404) = \qty{43}{mol}$,
$n_L = \qty{57}{mol}$.
\textbf{10.} 蒸气 $43 \times 0.626 = \qty{26.9}{mol}$；液体 $57 \times 0.404
= \qty{23.0}{mol}$；总计 \qty{49.9}{mol} $\approx 50$。
\textbf{11.} $26.9/50 = 54~\%$ 的苯，且蒸气中苯的含量仅为 63~\%。
\textbf{12.} 在 \qty{100}{\degreeCelsius} 时，平衡蒸气的 $y =
0.455 < 0.5$：进料点位于\omterm{def:b2:liquid-vapour-diagrams:bubble-dew}{露点线}上方，进料全部为
蒸气，因而无物质得以分离。
\textbf{13.} 在其\omterm{def:b2:liquid-vapour-diagrams:bubble-dew}{露点}处为 \qty{98.8}{\degreeCelsius}（方法同
\cref{exo:b2:liquid-vapour-diagrams:5}）。
\textbf{14.} $\alpha = p_1^*/p_2^*$；在 \qty{80.1}{\degreeCelsius} 时，$1.013/0.390 =
2.60$.
\textbf{15.} 在 \qty{110.6}{\degreeCelsius} 时，$2.377/1.013 = 2.35$。
\textbf{16.} $\sqrt{2.60 \times 2.35} = 2.47$。
\textbf{17.} $y = xp_1^*/p$ 与 $1 - y = (1 - x)p_2^*/p$；两式相除。
\textbf{18.} 无物质被引出：在相邻两块塔板之间，上升的蒸气与
下降的液体携带的物质的量以及各组分的量相同，
$V y_k = L x_{k+1}$，且 $V = L$，故 $x_{k+1} = y_k$。
\textbf{19.} 记 $r = x/(1 - x)$，则 $r(y_k) = \alpha\,r(x_k)$ 且 $x_{k+1} = y_k$，故
$r(x_{k+1}) = \alpha\,r(x_k)$ 及 $r(x_D) = \alpha^N r(x_B)$；$N = \ln[r(x_D)/
r(x_B)]/\ln\alpha$.
\textbf{20.} $N_{\min} = \ln(19 \times 19)/\ln 2.47 = 5.889/0.904 = 6.5$。
\textbf{21.} 精馏塔必须输出产品，因此它在有限的\omterm{def:b2:liquid-vapour-diagrams:fractional-distillation}{回流比}下运行，
这需要更多塔板；而实际塔板达不到平衡（其
效率低于 1）。
\textbf{22.} 共七步，最后一步结束于 0.034，低于 0.05：
按连续计数为 6.5 块塔板，按整数计数则为七块（再沸器计为一块）。
\textbf{23.} $x = (95/46.07)/(95/46.07 + 5/18.02) = 0.881$.
\textbf{24.} 塔顶至多能得到\omterm{def:b2:liquid-vapour-diagrams:azeotrope}{共沸物}（按摩尔计为 0.88）；塔底
得到水。五十块塔板也无法越过\omterm{def:b2:liquid-vapour-diagrams:azeotrope}{共沸物}。
\textbf{25.} 全回流时，$\boldsymbol{N_{\min} \approx 6.5}$ \omterm{def:b2:liquid-vapour-diagrams:fractional-distillation}{理论
塔板}。
\end{solution}
